\section{Omografia, IAC e Stima di \texorpdfstring{\(\mat{K}\)}{K}}
\label{sec:omografia_calibrazione}

%─────────────────────────────────────────────
\subsection{Coniche in \(\mathbb{P}^2\) e retroproiezione}
\label{sub:coniche}

Una conica in \(\mathbb{P}^2\) è descritta da una matrice
\(\mat{\Gamma}\) simmetrica \(3\times3\) con \(\det\mat{\Gamma}\neq0\)
tramite l'equazione
\[
  \homo{x}^T \mat{\Gamma}\, \homo{x} = 0.
\]
Se nello spazio 3D è presente una figura curva, i raggi visivi che
partono dal suo contorno e convergono in \(C\) formano un
\emph{cono 3D}; la sua intersezione con il piano del sensore produce
una \emph{conica 2D}. La \emph{retroproiezione} di \(\mat{\Gamma}\)
nello spazio 3D definisce la conica \(\mat{Q}\):
\[
  \mat{Q} = \mat{P}^T \mat{\Gamma}\, \mat{P},
  \qquad
  \homo{X}^T \mat{Q}\, \homo{X} = 0,
\]
dove \(\mat{Q}\) è \(4\times4\) simmetrica. Un punto \(\homo{X}_0\)
appartiene alla conica retroproiettata se e solo se la sua immagine
\(\mat{P}\homo{X}_0\) appartiene a \(\mat{\Gamma}\):
\[
  \homo{X}_0^T\, \mat{P}^T \mat{\Gamma}\, \mat{P}\,\homo{X}_0 = 0.
\]
Una proiettività \(\omega: \mathbb{P}^2 \to \mathbb{P}^2\) indotta da
\(\mat{A}\) con \(\det\mat{A}\neq0\) trasforma la conica secondo
\[
  \mat{\Gamma}' = (\mat{A}^{-1})^T \mat{\Gamma}\, \mat{A}^{-1},
\]
poiché sostituendo \(\homo{x} = \mat{A}^{-1}\homo{x}'\) nell'equazione
della conica si ottiene \(\homo{x}'^T\mat{\Gamma}'\homo{x}' = 0\).

%─────────────────────────────────────────────
\subsection{Conica assoluta \(\Omega\) e sua immagine}
\label{sub:conica_assoluta}

Una sfera nello spazio proiettivo ha equazione
\[
  X^2 + Y^2 + Z^2 + aX + bY + cZ + dU^2 = 0.
\]
Intersecandola con il piano all'infinito si ottiene la
\emph{conica assoluta} \(\Omega\):
\[
  \begin{cases}
    X^2 + Y^2 + Z^2 = 0 \\
    U = 0
  \end{cases}
\]
\(\Omega\) è indipendente dalla sfera scelta e non ha punti reali:
i suoi elementi sono i \emph{punti circolari impropri}
\((1:\pm i:0)\), comuni a qualsiasi piano \(\pi\) nello spazio.
Per qualsiasi piano \(\pi\), i suoi punti circolari impropri
\((1:\pm i:0)\in\pi\) si proiettano sul piano immagine tramite la
proiettività \(\mat{K}\mat{R}\) (che mappa il piano improprio
\(U=0\) sullo schermo). La conica immagine \(\mat{\Gamma}\) è:
\[
  \mat{\Gamma} = (\mat{K}\mat{R})^{-T}\,\mat{I}_3\,(\mat{K}\mat{R})^{-1}.
\]
Poiché \(\mat{R}\) è ortogonale (\(\mat{R}^{-1}=\mat{R}^T\)),
il termine rotazionale si semplifica:
\[
  \mat{\Gamma}
  = (\mat{K}^{-1})^T \underbrace{\mat{R}^{-T}\mat{R}^{-1}}_{=\,\mat{I}_3} \mat{K}^{-1}
  = (\mat{K}^{-1})^T \mat{K}^{-1}.
\]
\(\mat{R}\) scompare grazie all'ortogonalità: i valori misurati
sull'immagine dipendono quindi \emph{solo} da \(\mat{K}\).
L'inversa di \(\mat{\Gamma}\), nota come
\emph{Image of the Absolute Conic} (IAC), è:
\[
  \underbrace{
  \begin{bmatrix}
    f_1 & s   & x_0 \\
    0   & f_2 & y_0 \\
    0   & 0   & 1
  \end{bmatrix}
  }_{\mat{K}}
  \underbrace{
  \begin{bmatrix}
    f_1 & 0   & 0 \\
    s   & f_2 & 0 \\
    x_0 & y_0 & 1
  \end{bmatrix}
  }_{\mat{K}^T}
  =
  \underbrace{
  \begin{bmatrix}
    g^2 + x_0^2 & x_0 y_0     & x_0 \\
    x_0 y_0     & g^2 + y_0^2 & y_0 \\
    x_0         & y_0         & 1
  \end{bmatrix}
  }_{\mat{\Gamma}^{-1}}.
\]

\begin{figure}[H]
\centering
\begin{tikzpicture}[scale=1.2, >=Stealth]

  % ── Centro ottico C ──────────────────────────────────────────
  \coordinate (C) at (0, 0);
  \fill[black] (C) circle (2.5pt);
  \node[below left] at (C) {\(C\)};

  % ── Piano immagine (verticale a destra) ──────────────────────
  \draw[blue!70!black, thick]
    (3.5, -1.8) -- (3.5, 2.2);
  \node[blue!70!black, right] at (3.5, 2.0) {\small piano immagine};

  % ── Piano pi greco (inclinato, in basso a destra) ────────────
  \draw[orange!80!black, thick]
    (1.8, -2.0) -- (5.2, 0.6);
  \node[orange!80!black, right] at (5.2, 0.5) {\small \(\pi\)};

  % ── Cerchio sul piano pi greco (ellisse prospettica) ─────────
  \draw[orange!60!black, thick]
    (3.5, -0.9) ellipse (0.9cm and 0.28cm);

  % ── Punti circolari impropri I, J sul piano pi greco ─────────
  \coordinate (I)  at (2.62, -0.9);
  \coordinate (J)  at (4.38, -0.9);
  \fill[red!80!black] (I) circle (2.5pt);
  \fill[red!80!black] (J) circle (2.5pt);
  \node[below, font=\small, red!80!black] at (I) {\(I\)};
  \node[below, font=\small, red!80!black] at (J) {\(J\)};

  % ── Raggi da C verso I e J ────────────────────────────────────
  \draw[red!60!black, dashed, thick] (C) -- (I);
  \draw[red!60!black, dashed, thick] (C) -- (J);

  % ── Proiezioni di I e J sul piano immagine ───────────────────
  % Retta C->I interseca x=3.5:
  % parametrica: (0,0) + t*(2.62-0, -0.9-0), t = 3.5/2.62
  \coordinate (Iimg) at (3.5, -1.203);
  % Retta C->J: t = 3.5/4.38
  \coordinate (Jimg) at (3.5, -0.720);

  \fill[red!80!black] (Iimg) circle (2.5pt);
  \fill[red!80!black] (Jimg) circle (2.5pt);
  \node[right, font=\small, red!80!black] at (Iimg) {\(i\)};
  \node[right, font=\small, red!80!black] at (Jimg) {\(j\)};

  % ── Conica Gamma sul piano immagine ──────────────────────────
  \draw[blue!60!black, thick]
    (3.5, -0.96) ellipse (0.30cm and 0.52cm);
  \node[left, font=\small, blue!60!black] at (3.2, -0.96)
        {\(\mat{\Gamma}\)};

  % ── Asse ottico ───────────────────────────────────────────────
  \draw[->, gray!60, thin] (C) -- (4.5, 0)
        node[right, font=\small, gray!60] {\small asse ottico};

  % ── Etichetta KR ─────────────────────────────────────────────
  \node[font=\small, gray!40!black] at (1.6, 1.1)
        {\(\mat{K}\mat{R}\)};
  \draw[->, gray!50, thin] (1.6, 0.85) -- (2.4, 0.2);

\end{tikzpicture}
\caption{Proiezione dei punti circolari impropri \(I=(1:+i:0)\) e
         \(J=(1:-i:0)\) del piano \(\pi\) sul piano immagine tramite
         \(\mat{K}\mat{R}\). Le loro immagini \(i\) e \(j\) definiscono
         la conica \(\mat{\Gamma}\) (IAC) sul sensore.}
\label{fig:IAC_proiezione}
\end{figure}

%─────────────────────────────────────────────
\subsection{Stima di \(\mat{K}\) da omografie}
\label{sub:stima_K}

Il pattern usato nel progetto è planare (\(Z=0\)), quindi la terza
colonna di \(\mat{P}\) è ridondante e ogni vista \(j\) fornisce
un'\emph{omografia} \(\mat{M}_j\) (\(3\times3\), rango~3) che
correla i punti \((X,Y,0)\) del target con le loro proiezioni
\((u,v)\) sul sensore:
\[
  \varrho\,\homo{x} = \mat{M}_j\,\homo{X}_{2D},
  \qquad
  \homo{X}_{2D} = \begin{bmatrix} X \\ Y \\ 1 \end{bmatrix}.
\]
Denotando con \(\vect{h}_1, \vect{h}_2, \vect{h}_3\) le colonne
di \(\mat{M}_j\), l'ortogonalità di \(\mat{R}\) impone due vincoli:
\[
  \vect{h}_1 \cdot \mat{\Gamma}\, \vect{h}_2 = 0,
  \qquad
  \vect{h}_1 \cdot \mat{\Gamma}\, \vect{h}_1
  = \vect{h}_2 \cdot \mat{\Gamma}\, \vect{h}_2.
\]
Questi sono due equazioni lineari sui 6 elementi indipendenti di
\(\mat{\Gamma}\). Con almeno \(n\geq3\) viste il sistema è
determinato: si stima \(\mat{\Gamma}\), se ne ricava \(\mat{K}\)
per fattorizzazione di Cholesky su
\(\mat{\Gamma}^{-1} = \mat{K}\mat{K}^T\), e infine si ottengono
i parametri estrinseci \((\mat{R}_j, \vect{c}_j)\) per ogni frame
da
\[
  \mat{M}_j = \mat{K}
  \begin{bmatrix} \vect{r}_1 & \vect{r}_2 & \vect{t} \end{bmatrix}.
\]

%─────────────────────────────────────────────
\subsection{Modello di distorsione radiale e tangenziale}

Il modello pinhole suppone una trasformazione lineare, ma le lenti reali
introducono \emph{distorsione}: il punto effettivamente osservato
\((x_d,y_d)\) differisce dal punto teorico \((x_n,y_n)\) secondo
\begin{align}
  x_d &= x_n(1+k_1r^2+k_2r^4+k_3r^6)
        +2p_1x_ny_n+p_2(r^2+2x_n^2),\\
  y_d &= y_n(1+k_1r^2+k_2r^4+k_3r^6)
        +p_1(r^2+2y_n^2)+2p_2x_ny_n,
\end{align}
dove \(r^2=x_n^2+y_n^2\), \(k_1,k_2,k_3\) sono i coefficienti di
distorsione radiale (modello di Brown--Conrady) e
\(p_1,p_2\) quelli tangenziali.

I coefficienti \((k_1,k_2,k_3,p_1,p_2)\) non compaiono nella derivazione
lineare della IAC --- che è puramente proiettiva --- ma vengono stimati
congiuntamente a \(\mat{K}\) nella fase di raffinamento: una volta nota
una stima iniziale di \(\mat{K}\) tramite le omografie \(\mat{M}_j\),
l'ottimizzazione di Levenberg--Marquardt minimizza l'errore di
riproiezione tenendo conto dell'intera catena
\[
  (X,Y,Z) \;\xrightarrow{\mat{K},\,\mat{R}_j,\,\vect{c}_j}\;
  (x_n, y_n) \;\xrightarrow{k_1,k_2,k_3,\,p_1,p_2}\;
  (x_d, y_d) \;\xrightarrow{\text{pixel}}\; (u,v).
\]